% Source PDF: source/普通高考/2007/2007湖北理.pdf
% Page: 1
% PDF embedded bitmap attempt: pdfimages -list returned no image objects; the figure is vector content.
% Grid evidence: tmp/redraw_2007_hubei/page-1_grid100.png and tmp/redraw_2007_hubei/q17_block_grid50.png.
% Page crop evidence: tmp/redraw_2007_hubei/q17_orig.png (no-grid crop from the source page).
% Elements: horizontal axis labeled 纤度 with arrow, vertical axis labeled 频率/组距 with arrow, origin 0,
%   x-axis break mark, and ticks 1.30, 1.34, 1.38, 1.42, 1.46, 1.50, 1.54.
% Relations: the x-axis carries the fiber-fineness class boundaries; the y-axis is frequency density.
%   The source intentionally leaves the histogram plotting field blank for the student to complete.
% solid segments: x-axis, y-axis, arrowheads, tick marks, zigzag break
% dashed segments: none
% Label avoidance: axis labels are placed outside the plot; tick labels remain below the x-axis.
\begin{tikzpicture}
\begin{axis}[exam statistical histogram,
    width=10.0cm,
    height=4.2cm,
    xmin=0,
    xmax=3.3,
    ymin=0,
    ymax=1.05,
    axis lines=left,
    axis line style={->,black,line width=0.45pt},
    xtick={0,0.60,1.00,1.40,1.80,2.20,2.60,3.00},
    xticklabels={0,1.30,1.34,1.38,1.42,1.46,1.50,1.54},
    ytick=\empty,
    tick style={black,line width=0.4pt},
    tick align=outside,
    xlabel=\empty,
    ylabel=\empty,
    xticklabel style={font=\normalsize,anchor=north},
    clip=false,
]
\examstatylabel{\(\dfrac{\text{频率}}{\text{组距}}\)};
\node[anchor=south east,inner sep=0pt] at (axis cs:3.28,0.10) {纤度};
% The source has a broken x-axis between 0 and the first class boundary.
\draw[white,line width=2.0pt] (axis cs:0.18,0) -- (axis cs:0.40,0);
\examstatxbreak[black,line width=0.45pt]{0.29000000000000004}{0};
\end{axis}
\end{tikzpicture}
